% Appendix, from exam/README.md: what the two written papers are for, what they cover, why they ask
% for code, the conventions they assume, and how they are marked.

\chapter{The Self-Assessment Papers}
\label{exam:ch:about}

This appendix and the four after it hold two complete four-hour papers for \emph{Machine Learning},
with worked solutions. Appendices~\ref{exa:ch:paper} and \ref{exb:ch:paper} are the papers,
questions only; Appendices~\ref{ansa:ch:answers} and \ref{ansb:ch:answers} are their model answers,
with marks.

\section*{What these are for}
The book has \textbf{no exam and nothing is marked}. What you get instead is a test suite for the
code exercises of Chapters~1 to~5 and~8 to~10, the conv demo of Chapter~\ref{c6:ch:cnn} printing its
results next to the numbers you worked out by hand in that chapter's hand-training exercise, the
worked solutions in Appendix~\ref{sol:ch:solutions}, and the self-assessment exercises in
\secref{c10:sec:exercises}. Every one of those checks a component you built or a number you computed
beside it.

These papers do not replace that. \textbf{They are here to test your own skills and knowledge,
nothing more.} They gate nothing, they are not a qualification, and no part of the book requires
them. Nothing in the companion repository depends on them and \code{make test} does not know they
exist. They are useful where a written result is wanted anyway, for a certifying employer or a
formal course credit, and useful on their own for finding out what you can reconstruct with nothing
in front of you.

\textbf{Take one after the last chapter.} Every paper draws on all ten chapters, so sitting one
partway through examines material you have not read yet, and the result says more about how far you
have read than about what you have understood. The intended point is after
Chapter~\ref{c10:ch:flatten}, once your \code{cnn_work} trains end to end.

\section*{The two papers}
Both cover the whole book, Chapters~1 to~10, and share no question. Either can be used alone; use
both as a main sitting and a resit, or in alternate years.

They do \textbf{not} weight the book identically, and that is deliberate rather than accidental.
Paper~A spends more on the convolutional layer itself (Chapter~6 is worth roughly three times as
much there), Paper~B spends more on the finished network and its constraints (Chapter~10 is worth
twice as much). Two topics appear on one paper only: the precision and early stopping of
Chapter~\ref{c2:ch:training} are examined on \textbf{B alone} (Q2c, Q2d), and the fixed-point and
real-time material from \secref{c10:sec:embedded} is likewise \textbf{B alone} (Q9b, Q9c). Paper~A
is also the only one that asks a candidate to write a method from scratch. If you are using the
papers as a sitting and a resit, that asymmetry is worth knowing; if you need the two to examine
exactly the same ground, sit both.

\subsection*{Paper A}

\begin{booktable}{@{}p{4.6em}Lp{5.4em}r@{}}
  \thead{Question} & \thead{Topic} & \thead{Chapters} & \thead{Marks} \\
  \midrule
  1 & Learning the rules instead of writing them & 1, 2 & 12 \\
  2 & The training loop that trains nothing & 2, 4 & 12 \\
  3 & A network trained by hand & 3 & 14 \\
  4 & The activation function, and the argument it is given & 3, 5 & 10 \\
  5 & The dense layer, written out & 5 & 12 \\
  6 & Why an image is not just a long vector & 6 & 10 \\
  7 & A convolutional layer, traced by hand & 6, 7, 8 & 14 \\
  8 & Max pooling, and the value it forgets & 7, 9 & 8 \\
  9 & The finished network, and the target it has to run on & 10 & 8 \\
  \midrule
    & & & \textbf{100} \\
\end{booktable}

\subsection*{Paper B}

\begin{booktable}{@{}p{4.6em}Lp{5.4em}r@{}}
  \thead{Question} & \thead{Topic} & \thead{Chapters} & \thead{Marks} \\
  \midrule
  1 & The model, and the rate that destroys it & 1 & 10 \\
  2 & Order, precision, and knowing when to stop & 2, 4 & 10 \\
  3 & Two output nodes & 3 & 14 \\
  4 & The network class, and what a stub can prove & 3, 4 & 12 \\
  5 & The dense layer, and the mistake only Tanh can catch & 5 & 12 \\
  6 & Padding, sharing, and the kernel's gradient & 6, 7, 8 & 14 \\
  7 & A pooling layer that runs and returns true & 7, 9 & 12 \\
  8 & Every layer, side by side & 10 & 8 \\
  9 & Off the laptop & 10 & 8 \\
  \midrule
    & & & \textbf{100} \\
\end{booktable}

Paper~A leans towards \textbf{tracing an algorithm forwards}: one epoch of linear regression by hand,
a full \code{feedforward()}/\code{backpropagate()}/\code{optimize()} step through a small network, a
$4 \times 4$ image pushed through conv, pooling and flatten, and the gradients routed back.
Paper~B leans towards \textbf{the conditions under which the algorithm stops working}: the learning
rate that diverges inside the range \code{train()} permits, the derivative that is right for ReLU and
wrong for Tanh, the even kernel whose padding does not line up, and a pooling layer whose faults never
produce an error message.

\section*{Code on a written paper}
Unlike the papers for some other QAcademy courses, these two \textbf{do ask a candidate to write
C++}. That is deliberate. This book teaches machine learning by building it: the specification of
every layer is prose in a chapter and a test suite that tells you whether what you wrote satisfies
it, and a paper that avoided code entirely would be examining a different book.

The code questions come in two shapes, and neither of them is a typing exercise:
\begin{itemize}
  \item \textbf{Write one method.} \code{Dense::feedforward()}, \code{backpropagate()} and
    \code{optimize()} are asked for in Paper~A; the answers are ten to twenty lines each. What
    carries the marks is the algorithm, namely the validation, the weighted sum starting at the
    bias, the derivative taken at the right argument; not the syntax.
  \item \textbf{Find what is wrong with the algorithm.} Both papers hand over a listing that
    compiles, runs, and returns \code{true}, and ask what it does instead of what it should. The
    faults are the ones the course's own test-suite READMEs single out: a shuffle nobody reads,
    backpropagation in the wrong order, a max search that starts at \code{0.0}, a gradient matrix
    that is never cleared.
\end{itemize}
Both papers state in their rubric that answers may be written in C++17 \textbf{or in unambiguous
pseudocode}, and that nothing is marked on syntax. A missing semicolon costs nothing. A missing
bounds check costs everything.

\section*{Conventions the papers assume}
Both papers state these in their own rubric, so a candidate never has to have read this appendix.
\begin{itemize}
  \item \textbf{The book's sign convention.} The deviation is $\delta = \yref - \yp$ and every
    parameter is \emph{increased} by its adjustment, so every update reads
    \code{parameter += gradient * learningRate} with no minus sign. A candidate who flips the sign
    consistently and says so is not penalized twice.
  \item \textbf{The activation derivative is evaluated at the weighted sum $s$}, never at the output
    $y$. Both papers have a question about what happens when it is not.
  \item \textbf{Conv layers zero-pad} with \code{pad = kernelSize / 2}, so the output is the same
    size as the input.
  \item \textbf{Max pooling breaks ties toward the first occurrence} in row-major order, and
    \textbf{flatten is row-major}, matching the hand-training exercise of
    Chapter~\ref{c6:ch:cnn} and the shipped implementations.
  \item \textbf{Only the values nobody could be expected to carry are supplied}: a handful of
    \code{tanh()} values, $2^{16}$ on Paper~B, and that a KB is 1024 bytes. Parameter counts, output
    shapes and padding sizes are derived, because deriving them is the examinable skill.
  \item \textbf{Code questions refer to the \code{Dense}/\code{Shallow} API specified in
    Chapters~3 to~5}, the one the candidate builds, not to the refactored API in \code{cnn_demo}.
    The two differ: \code{cnn_demo}'s \code{Dense} reverses the constructor arguments, renames the
    accessors, collapses the two \code{backpropagate()} overloads into one, and carries a seventh
    member variable. Answers written against either are accepted, but the model answers follow
    Chapters~3 to~5.
\end{itemize}

\section*{Marking}
Every solution is written to be marked by somebody who has read the chapters and does not otherwise
write machine learning code, so each carries the reasoning rather than the answer alone. Marks are
shown per part in the solutions.

Three conventions worth agreeing before a paper is marked:
\begin{itemize}
  \item \textbf{Method carries the marks.} A correct update rule with an arithmetic slip in it is
    worth more than a correct number with no working, and both papers are built so that later parts
    consume earlier answers. Follow through an error rather than penalizing it twice.
  \item \textbf{The named traps are worth full marks on their own.} Several questions exist entirely
    to see whether a candidate avoids a specific mistake: the activation derivative taken from $y$
    instead of $s$, a gradient that ReLU has already annihilated, \code{optimize()} given the
    network's input instead of the previous layer's output, a max search starting at \code{0.0}, an
    \code{optimize()} invented for a layer that has nothing to optimize. Where a solution flags one
    of these, a candidate who walks into it loses those marks and no others.
  \item \textbf{The discussion parts are not decoration.} ``State why'' is where the book's actual
    content is, and a paper marked only on the arithmetic would pass a candidate who has understood
    none of it.
\end{itemize}
